\chapter{CONCLUSÃO}\label{chap:conclusao}

O protótipo demonstra os elementos centrais de comunicação de grupo e ordenação
total em um chat distribuído. Cada participante possui processo, porta e estado
próprios; mensagens de grupo atravessam sockets TCP e recebem uma sequência única
do coordenador. O buffer por lacunas faz com que mensagens recebidas fora de ordem
sejam entregues somente quando todas as posições anteriores estão disponíveis.
Nos ensaios com 3, 8 e 15 nós, todos os participantes terminaram com históricos
globais idênticos e buffers vazios.

O relógio vetorial acompanha as mensagens e torna visível a evolução causal, mas
a ordem total é decidida pelo sequenciador, como previsto pela abordagem B. O
monitoramento por heartbeat e a eleição do Valentão conseguiram selecionar o maior
ID ativo após a queda do líder em um cenário com três nós. A arquitetura, portanto,
oferece uma base concreta para estudar comunicação, relógios lógicos, difusão e
coordenação. O comando \texttt{snapshot} acrescenta uma fotografia consistente dos
estados locais e canais sem suspender a aplicação. Nos ensaios, a captura agregou
3 e 15 relatórios, e o cenário controlado preservou corretamente uma mensagem em
trânsito.

O R6 é atendido pela implementação de Chandy--Lamport, apoiada pela numeração FIFO
dos canais lógicos e pela agregação automática dos fragmentos no iniciador. A
continuidade da ordem total também foi demonstrada: depois da queda do coordenador
inicial, nó 3, o nó 2 recuperou o prefixo global, adotou a próxima sequência
correta e entregou o pedido que havia detectado a falha. A comunicação privada foi
validada separadamente e não interfere na ordem global das mensagens de grupo.

Os requisitos funcionais do roteiro foram demonstrados nos cenários executados.
Como evolução, recomenda-se reduzir o custo de retransmitir todo o histórico na
recuperação, adicionar persistência opcional e tratar a falha de um participante
durante um snapshot com timeout e nova configuração de membros. Essas limitações
não alteraram a convergência nos testes apresentados, mas delimitam o comportamento
do protótipo sob falhas sucessivas ou execuções prolongadas.
